\begin{helper}
Fix a history cell with recorded coordinate set \(R_h\), terminal coordinate
set \(U_h\), and two completions \(\omega,\omega^\ast\) having the same
injection.  Assume that both samples have the same pre-terminal recorded set:
\[
e\in R_h
\quad\Longleftrightarrow\quad
e\in\noderef{TwoBitePreTerminalRecordedPairs}(\omega,\varepsilon,I)
\quad\Longleftrightarrow\quad
e\in\noderef{TwoBitePreTerminalRecordedPairs}(\omega^\ast,\varepsilon,I),
\]
and that the terminal set is the full unrecorded terminal universe:
\[
e\in U_h
\quad\Longleftrightarrow\quad
e\in\noderef{TwoBiteTerminalCoordinateUniverse}(m)\hbox{ and }e\notin R_h.
\]
Assume also that \(\omega\) and \(\omega^\ast\) agree on every coordinate of
\(R_h\).

If they agree on all blue coordinates in \(U_h\), then for every blue base
vertex \(b\) the sets
\[
\noderef{TwoBiteX}(\omega,I,\operatorname{blue}(b))
\quad\text{and}\quad
\noderef{TwoBiteX}(\omega^\ast,I,\operatorname{blue}(b))
\]
are equal.  Consequently the large, medium, and small classifications of
\(\operatorname{blue}(b)\) agree in the two samples, and for every base pair
\(q\),
\[
q\in\noderef{TwoBiteBasePairs}
\bigl(\pi_R^\omega(\noderef{TwoBiteX}(\omega,I,\operatorname{blue}(b)))\bigr)
\]
if and only if
\[
q\in\noderef{TwoBiteBasePairs}
\bigl(\pi_R^{\omega^\ast}
(\noderef{TwoBiteX}(\omega^\ast,I,\operatorname{blue}(b)))\bigr).
\]
The color-swapped statement also holds: agreement on all red coordinates in
\(U_h\) transfers the red sets
\(\noderef{TwoBiteX}(-,I,\operatorname{red}(r))\), their large/medium/small
classifications, and the corresponding blue projected-pair predicates.
\end{helper}
\begin{proof}
We prove the blue-neighborhood statement; the red statement is identical with
the color tags interchanged.  Fix a blue base vertex \(b\) and a final vertex
\(v\in I\).  Membership of \(v\) in
\(\noderef{TwoBiteX}(\omega,I,\operatorname{blue}(b))\) is, by unfolding
\(\noderef{TwoBiteX}\) and \(\noderef{TwoBiteLiftedNeighborhood}\), the
assertion that the blue base edge between \(b\) and
\(\pi_B^\omega(v)\) is present, except in the loop case where the simple blue
base graph has no loop.  The common-injection hypothesis gives
\[
\pi_B^\omega(v)=\pi_B^{\omega^\ast}(v).
\tag{1}
\]
Thus the same base endpoint is tested in both samples.

If this endpoint is \(b\), the queried adjacency is a loop in both blue base
graphs and is false in both samples.  Hence the membership decision agrees.
Assume now that the endpoint is distinct from \(b\).  Put the two endpoints
in increasing order and call the resulting tagged coordinate \(c\).  Then
\[
c=\operatorname{blue}(a_1,a_2),\qquad a_1<a_2,
\]
so by the definition of \(\noderef{TwoBiteTerminalCoordinateUniverse}\) we
have \(c\in\noderef{TwoBiteTerminalCoordinateUniverse}(m)\).  There are two
cases.

First, if \(c\in R_h\), then the recorded-agreement hypothesis gives
\[
\noderef{TwoBiteEdgeCoordinateValue}(\omega,c)
\Longleftrightarrow
\noderef{TwoBiteEdgeCoordinateValue}(\omega^\ast,c).
\tag{2}
\]
Since \(c\) is the blue coordinate representing the unordered pair
\(\{b,\pi_B(v)\}\), (2) is exactly equality of the blue adjacency truth value
used in the two lifted-neighborhood tests.

Second, if \(c\notin R_h\), then the full-terminal-universe identity gives
\[
c\in U_h.
\tag{3}
\]
The assumed agreement on blue terminal coordinates applies to (3), and again
gives the equivalence of the blue edge value in the two samples.  These two
cases cover every non-loop blue coordinate because every such coordinate lies
in the terminal coordinate universe and \(U_h\) is precisely the complement
of \(R_h\) inside that universe.

We have proved that every \(v\in I\) is a blue neighbor of \(b\) in
\(\omega\) if and only if it is a blue neighbor of \(b\) in
\(\omega^\ast\).  Since
\(\noderef{TwoBiteX}\) is the filter of \(I\) by this lifted-neighborhood
predicate, the two finite sets \(X_b(I)\) are equal.  Their cardinalities are
therefore equal.  Unfolding
\(\noderef{TwoBiteLargeClass}\), \(\noderef{TwoBiteMediumClass}\), and
\(\noderef{TwoBiteSmallClass}\), each of the three predicates is a numerical
condition only on \(|X_b(I)|\) and the fixed cutoffs.  Hence the three
classifications agree between \(\omega\) and \(\omega^\ast\).

It remains to transfer the projected-pair predicate.  The common injection
also gives
\[
\pi_R^\omega(v)=\pi_R^{\omega^\ast}(v)
\qquad(v\in I).
\tag{4}
\]
Together with the equality of \(X_b(I)\), equation (4) gives equality of the
red projected images
\[
\pi_R^\omega\bigl(\noderef{TwoBiteX}(\omega,I,\operatorname{blue}(b))\bigr)
=
\pi_R^{\omega^\ast}
\bigl(\noderef{TwoBiteX}(\omega^\ast,I,\operatorname{blue}(b))\bigr).
\]
Applying the same finite-set function
\(\noderef{TwoBiteBasePairs}\) to the two equal projected images gives the
displayed equivalence for every base pair \(q\).

For the red statement, replace blue by red throughout.  A vertex
\(v\in I\) lies in
\(\noderef{TwoBiteX}(-,I,\operatorname{red}(r))\) according to the red edge
coordinate between \(r\) and \(\pi_R(v)\).  That coordinate is either in the
recorded set \(R_h\), where recorded agreement applies, or in the terminal
set \(U_h\), where the assumed red terminal agreement applies.  The same
cardinality and projection-image argument transfers the red
large/medium/small classifications and the blue projected-pair predicates.
\end{proof}
